% 级数（综述）
% license CCBYSA3
% type Wiki

本文根据 CC-BY-SA 协议转载翻译自维基百科\href{https://en.wikipedia.org/wiki/Series_(mathematics)}{相关文章}

在数学中，级数大致而言就是将无限多个项一个接一个地相加。[1]对级数的研究是微积分及其推广——数学分析——的重要组成部分。级数在数学的大多数领域都有应用，甚至在组合学中研究有限结构时，也会通过生成函数来使用级数。无限级数的数学性质使其在物理学、计算机科学、统计学和金融等其他定量学科中也得到了广泛应用。

在古希腊人看来，一个“潜在无限”的求和可以得到一个有限的结果，这被认为是悖论，其中最著名的例子就是芝诺悖论。[2][3]尽管如此，古希腊数学家们，包括阿基米德，仍然在实践中使用了无限级数，例如在抛物线求积问题中。[4][5]芝诺悖论中的数学问题在 17 世纪通过“极限”概念得到了化解，尤其是通过艾萨克·牛顿的早期微积分。[6]在19世纪，通过卡尔·弗里德里希·高斯和奥古斯丁-路易·柯西等人的工作，这一解决方案得到了进一步的严格化与改进。[7]他们回答了哪些和式是存在的问题（通过实数的完备性），以及级数项是否可以重新排列而不改变其和的问题（通过绝对收敛与条件收敛的概念）。

在现代术语中，任何一个有序的无限序列$(a_{1}, a_{2}, a_{3}, \ldots)$，无论这些项是数字、函数、矩阵，还是其他可以相加的对象，都定义了一个级数，即把这些$a_i$ 一个接一个地相加。为了强调项的数量是无限的，级数常常也被称为无限级数，以区别于有限级数，后者有时用来指有限和。级数通常写作：
$$
a_{1} + a_{2} + a_{3} + \cdots ,~
$$
或者使用大写希腊字母 Σ 的求和符号表示：[8]
$$
\sum_{i=1}^{\infty} a_i .~
$$
由于级数所表达的无限次加法无法在有限时间内逐一完成，但如果这些项及其有限部分和属于一个具有极限的集合，就可能为该级数赋予一个值，称为级数的和。这个值定义为：当 $n \to \infty$ 时，级数前$n$项的部分和的极限（若该极限存在）。[9][10][11] 这些有限和称为部分和。用求和符号表示：
$$
\sum_{i=1}^{\infty} a_i = \lim_{n \to \infty} \sum_{i=1}^{n} a_i ,~
$$
若极限存在。[9][10][11]当极限存在时，该级数称为收敛或可和，此时序列$(a_{1}, a_{2}, a_{3}, \ldots)$ 也是可和的；反之，如果极限不存在，则称该级数发散。[9][10][11]
表达式$\sum_{i=1}^{\infty} a_{i}$既表示级数——即无限地一个接一个相加项的隐含过程——又在该级数收敛时表示级数的和——即这一过程的明确极限。这是对类似约定的一种推广：例如$a+b$既表示“加法”这一过程，也表示其结果——$a$ 与 $b$ 的和。

通常，级数的项来自某个环，最常见的是实数域$\mathbb{R}$或复数域$\mathbb{C}$。在这种情况下，所有级数的集合本身也构成一个环，其中的加法是逐项相加，而乘法则由柯西乘积定义。[12][13][14]
\subsection{定义}
\subsubsection{级数}
一个级数，或更完整地说，无限级数，就是一个无限和。它通常表示为：[8][15][16]
$$
a_{0} + a_{1} + a_{2} + \cdots \quad \text{或} \quad a_{1} + a_{2} + a_{3} + \cdots ,~
$$
其中各项 $a_{k}$ 是一个数列（序列）的成员，这些成员可以是数、函数，或任何可以相加的对象。级数也常用大写希腊字母 Σ 的符号来表示：[8][16]
$$
\sum_{k=0}^{\infty} a_{k} \quad \text{或} \quad \sum_{k=1}^{\infty} a_{k}.~
$$
另外，一种常见的表示方法是写出前若干项，加上省略号，再写出一般项，最后再加省略号，其中一般项是$n$的函数形式的第$n$项：
$$
a_{0} + a_{1} + a_{2} + \cdots + a_{n} + \cdots 
\quad \text{或} \quad 
f(0) + f(1) + f(2) + \cdots + f(n) + \cdots .~
$$
例如，欧拉常数$e$可以通过以下级数定义：
$$
\sum_{n=0}^{\infty} \frac{1}{n!} 
= 1 + 1 + \frac{1}{2} + \frac{1}{6} + \cdots + \frac{1}{n!} + \cdots ,~
$$
其中$n!$表示前$n$个正整数的乘积，而按照约定$0! = 1$。[17][18][19]
\subsubsection{级数的部分和}
给定一个级数$s = \sum_{k=0}^{\infty} a_k $,它的第 $n$ 个**部分和**定义为：[9][10][11][16]
$$
s_{n} = \sum_{k=0}^{n} a_k = a_0 + a_1 + \cdots + a_n .~
$$
有些作者会直接把一个级数与其部分和序列等同。[9][11] 无论是部分和序列还是项的序列，都能完全刻画这个级数。而且，项的序列可以通过部分和序列的相邻差来恢复：
$$
a_{n} = s_{n} - s_{n-1}.~
$$
数列的部分求和是线性序列变换的一个例子，在计算机科学中也称为前缀和。而从部分和恢复原数列的逆变换就是有限差分，这也是另一种线性序列变换。

有时，级数的部分和具有更简单的闭合形式。例如：等差级数的部分和：
$$
s_{n} = \sum_{k=0}^{n} (a + kd) 
= a + (a+d) + (a+2d) + \cdots + (a+nd) 
= (n+1)\left(a + \tfrac{1}{2}nd\right) .~
$$
等比级数的部分和：[20][21][22]
$$
s_{n} = \sum_{k=0}^{n} ar^{k} 
= a + ar + ar^{2} + \cdots + ar^{n} 
= a \cdot \frac{1-r^{\,n+1}}{1-r}, \quad \text{当 } r \neq 1,
$$

或者更简单地写成：

$$
s_{n} = a(n+1), \quad \text{当 } r = 1.
$$

---

要不要我接下来帮你翻译“级数的收敛与发散”这一节？
